\chapter{Implementation}

\section{System Architecture}

The experimental framework was implemented as an automated pipeline designed to analyze, modify, and evaluate open-source Python repositories. The system follows a modular architecture based on the Pipes and Filters architectural pattern, where each processing stage is implemented as an independent filter component.

This design enables reproducibility, extensibility, and isolation of experimental stages. Each filter receives structured input, performs a specific task, and produces output that is passed to the next stage.

The overall implementation consists of the following main subsystems:

\begin{itemize}
\item Dataset preparation subsystem
\item Static analysis subsystem
\item LLM-based architecture improvement subsystem
\item Maintainability evaluation subsystem
\item Results storage and artifact management subsystem
\end{itemize}

The entire pipeline operates automatically without manual intervention.

\section{Technology Stack}

The system was implemented using Python due to its extensive ecosystem, availability of static analysis tools, and compatibility with modern Large Language Models.

The core technologies used in the implementation are summarized below:

\begin{itemize}
\item Programming language: Python 3.12
\item Version control interaction: Git
\item Repository hosting platform: GitHub
\item Static analysis tool: Radon
\item LLM access: Ollama API
\item Data storage format: CSV and JSON
\item Architecture pattern: Pipes and Filters
\end{itemize}

Python was selected due to its suitability for automation, availability of maintainability analysis tools, and relevance to backend software development.

\section{Dataset Preparation Implementation}

The dataset preparation process was implemented as a sequence of automated filters.

The first component collected candidate repositories from GitHub using predefined search criteria. Repository metadata, including name, clone URL, and star count, was retrieved via the GitHub REST API.

The second component cloned each repository locally using shallow cloning to minimize network and storage overhead:

\begin{verbatim}
git clone --depth 1 <repository_url>
\end{verbatim}

Each repository was analyzed to determine whether it represents a backend system. The following conditions were verified:

\begin{itemize}
\item Presence of a \texttt{requirements.txt} file
\item Presence of Python source code files
\item Absence of machine learning libraries in dependencies
\item Absence of experimental or notebook-based project structures
\end{itemize}

Repositories that satisfied these criteria were included in the dataset.

Repositories were deleted immediately after analysis to minimize storage usage.

\section{Static Analysis Implementation}

Maintainability metrics were collected using automated static analysis.

The Radon tool was used to compute Maintainability Index values:

\begin{verbatim}
radon mi --json <repository_path>
\end{verbatim}

The output was parsed and stored in structured JSON format.

Three categories of maintainability metrics were collected:

\begin{itemize}
\item Code maintainability metrics
\item Architecture metrics
\item Documentation metrics
\end{itemize}

Results were stored in the artifacts directory structure:

\begin{verbatim}
artifacts/static_analysis/BEFORE/<repository_name>/
artifacts/static_analysis/AFTER/<repository_name>/
\end{verbatim}

This enabled direct comparison between the original and modified repository versions.

\section{LLM Integration Implementation}

Large Language Models were integrated to perform architectural analysis and improvements.

The LLM was used for three independent tasks:

\begin{itemize}
\item Architecture detection
\item Architectural tactic selection
\item Architectural tactic implementation
\end{itemize}

Repository structure and file contents were extracted and provided to the model as structured input.

To reduce token usage, only relevant information was included:

\begin{itemize}
\item Repository file tree
\item File contents (imports and signatures)
\item Static analysis results
\end{itemize}

LLM responses were required to follow strict JSON format, which enabled automatic parsing and execution.

\section{Architectural Tactic Implementation}

Architectural improvements were implemented automatically using LLM-generated code modifications.

The system enforced Test-Driven Development principles during implementation.

Each modification followed the sequence:

\begin{enumerate}
\item Generate a failing test
\item Implement minimal production code
\item Verify test execution
\end{enumerate}

The LLM generated structured instructions specifying file creation or modification.

These instructions were executed automatically within the repository.

All modifications were stored locally.

\section{Maintainability Re-evaluation}

After architectural improvements were applied, static analysis was executed again.

The same metrics collected before modification were recalculated for the modified repository.

This ensured consistency and comparability of results.

The Maintainability Index average value was used as the primary quantitative indicator.

\section{Data Storage and Artifact Management}

All experimental results were stored for further analysis and reproducibility.

The system stored:

\begin{itemize}
\item Repository metadata
\item Static analysis results
\item LLM outputs
\item Maintainability metrics
\item Final evaluation dataset
\end{itemize}

The final dataset was stored in CSV format.

Each record contained:

\begin{itemize}
\item Repository name
\item Repository URL
\item Stars count
\item Maintainability Index before improvement
\item Maintainability Index after improvement
\end{itemize}

This dataset was used for final evaluation.

\section{Automation and Execution Workflow}

The entire experimental process was automated.

For each repository, the system performed the following steps:

\begin{enumerate}
\item Clone repository
\item Run static analysis
\item Detect architecture using LLM
\item Select architectural tactic
\item Implement tactic
\item Run static analysis again
\item Store results
\item Delete repository
\end{enumerate}

This process was repeated sequentially for all repositories in the dataset.

Sequential execution ensured isolation and prevented interference between experiments.

\section{Reproducibility Considerations}

To ensure reproducibility, the implementation incorporated the following measures:

\begin{itemize}
\item Fixed analysis procedure
\item Automated execution
\item Structured artifact storage
\item Version-controlled codebase
\item Deterministic static analysis tools
\end{itemize}

All experimental results can be reproduced by executing the pipeline with the same dataset.
